%% kapa: exam=1 round=27 year=2016 period=1 form=A subject=civil_law subject_ko=민법 seq=3 number=3
%% source: Q-Net 감정평가사 기출문제 (cst00309) artlSeq=5169402 file=extracted/1008/5169402/제27회 감정평가사 1차 문제/제27회 감정평가사 1교시문제(A형).pdf
%% conversion: kordoc 4.24.51 -> tools/gen1.py (auto); review=proofread-sonnet; flags=none
\begin{question}{3}
사단법인의 사원총회에 관한 설명으로 옳지 않은 것은?
\begin{choices}
  \item 사원총회에는 대외적인 대표권이나 대내적인 업무집행권이 없다.
  \item 각 사원은 평등한 결의권을 가지며, 정관으로도 달리 정할 수 없다.
  \item 정관에 다른 규정이 없는 한, 총사원의 5분의 1 이상이 회의의 목적사항을 제시하여 총회 소집을 청구한 경우에 이사는 임시총회를 소집하여야 한다.
  \item 총회는, 정관에 규정이 있으면, 소집 통지에 기재한 목적사항 이외에 대해서도 결의할 수 있다.
  \item 정관에 다른 규정이 없는 한, 정관변경을 위해서는 총사원의 3분의 2 이상의 동의가 있어야 한다.
\end{choices}
\end{question}
\answer{2}
